    % University computer vision project in the CVPR 2026 format.
% The official CVPR style is used in final mode so author names remain visible.
% Line numbers are enabled separately for the course submission.
\documentclass[10pt,twocolumn,letterpaper]{article}

\usepackage{cvpr}
\usepackage[switch,mathlines]{lineno}

\definecolor{cvprblue}{rgb}{0.21,0.49,0.74}
\usepackage[pagebackref,breaklinks,colorlinks,allcolors=cvprblue]{hyperref}

% Replace each visible TODO as the project develops.
\newcommand{\papertodo}[1]{\textit{[TODO: #1]}}

\title{How Can Deep Learning Improve [Computer Vision Task]?\\[0.4em]
{\large Deep Learning for Computer Vision Project}}
\author{Erdem Bayrakcioglu \qquad Henricus Bracht \qquad Leonard Fischbach\\
Universit\`a Bocconi}

\begin{document}
\maketitle
\linenumbers

\begin{abstract}
\papertodo{Summarize the problem, approach, main results, and conclusion.}
\end{abstract}

\section{Introduction}
\papertodo{Introduce the computer vision problem, its motivation, and the project contributions.}

\section{Related Work}
\papertodo{Discuss the most relevant prior methods and explain how this project relates to them. Add citations.}

\subsection{Vision-Language-Action Models}
\papertodo{Short explanation of how models translate images and language instruction into roboter actions. RT-2 and Open-VLA as start, then $\pi_0$ as the introduction of Movement-Action exper with flow matching.}

\subsection{$\pi_0.5$}
\papertodo{Explain shortly what $\pi_0.5$ has contributed and whats new with it compared to former VLAs}

\subsection{Xiaomi Robotics 1}
\papertodo{Explain shortly what XR1 has contributed and whats new with it compared to former VLAs, especially whats different from $pi_0.5$}

\subsection{Robustness and Distribution Shifts}
\papertodo{IMPORTANT: Explain why a high success rate under ideal conditions is not 100\% meaningful under chaging conditions}

\subsection{Research Gap and Our Contribution}
\papertodo{What does our investigation deliver? -> A controlled comparison of the two models with the same task the same visual disturbances and an analysis of the errors.}

\section{Method}
\papertodo{Describe the dataset, preprocessing, model or algorithm, and implementation details.}

\section{Experimental Setup}
\papertodo{State the data splits, baselines, evaluation metrics, and training or inference settings.}

\section{Results and Discussion}
\papertodo{Present quantitative results, qualitative examples, comparisons, and interpretation. Add figures and tables as needed.}

\section{Conclusion}
\papertodo{Summarize the findings, limitations, and possible future work.}

% Enable the bibliography after adding references and citations:
% {\small
% \bibliographystyle{ieeenat_fullname}
% \bibliography{references}
% }

\end{document}
